\documentclass[journal]{IEEEtran}
\usepackage{amsmath, amssymb, graphicx, booktabs, hyperref}
\usepackage{cite}

\title{PHANTOM: Kinematically Consistent False LOS Rate Injection
for Controlled Miss Distance via EKF Innovation Gate Exploitation}

\author{Naresh~Dama%
\thanks{N.~Dama is with the Department of Aerospace and Mechanical
Engineering, Jacksonville, FL 32256 USA
(e-mail: naresh.dama@unf.edu).}%
\thanks{Manuscript submitted March 2026.}}

\begin{document}
\maketitle

\begin{abstract}
Modern radar-guided missiles employ Extended Kalman Filters (EKF)
with chi-squared innovation gating to achieve robust target tracking
against noisy measurements.  This paper presents PHANTOM
(Proportional-Navigation Heuristic Adaptive Missile-Tracking
Nullification), a simulation framework demonstrating that
kinematically consistent false line-of-sight (LOS) rate injection
can systematically exploit the EKF gating mechanism to produce
controlled miss distances while evading statistical detection.
A linear ramp injection profile at $0.035$~rad/s achieves a mean
miss distance of $20\,536 \pm 909$~m across $N=200$ Monte~Carlo
engagements, with an EKF detection rate of $4.22\%$ --- below the
$5\%$ operational threshold.  The difference in miss distance
between baseline and injection engagements is statistically
significant ($t = 104.36$, $p < 0.001$, Cohen's~$d = 58.16$).
Miss distance scales as a power law with exponent $b = 1.53$
($R^2 = 0.940$) over the pre-saturation ramp-rate regime,
consistent with quadratic false-angle accumulation.  All code
and data are publicly archived at \texttt{10.5281/zenodo.XXXXXX}.
\end{abstract}

\begin{IEEEkeywords}
missile guidance, Extended Kalman Filter, proportional navigation,
electronic warfare, false target injection, chi-squared gating,
Monte~Carlo simulation
\end{IEEEkeywords}

%% ================================================================
\section{Introduction}
\label{sec:intro}

Proportional Navigation (PN) guidance, combined with an Extended
Kalman Filter (EKF) for target tracking, forms the backbone of
modern radar-guided missile seekers.  When the EKF operates
correctly, the system achieves sub-metre miss distances against
maneuvering targets, yielding kill probabilities exceeding
$0.9$~\cite{zarchan2012}.  The high lethality of PN-guided
threats motivates the electronic warfare (EW) community to develop
countermeasures that degrade the seeker's tracking accuracy.

Conventional countermeasures---chaff, jamming, and Digital
Radio-Frequency Memory (DRFM) repeaters---generally aim to
saturate, blind, or confuse the radar receiver.  These approaches
often trigger the very anomaly-detection mechanisms (innovation
gating, adaptive thresholding) designed to reject spurious
measurements, leading to rapid identification and counter-counter
measures~\cite{barshalom2001}.

PHANTOM proposes an alternative philosophy: instead of overwhelming
the EKF, \emph{operate within} its acceptance envelope.  By
injecting a slowly ramping false LOS rate whose instantaneous
contribution to the innovation statistic remains below the
chi-squared gate, the injection corrupts the target-state estimate
cumulatively without triggering measurement rejection.  The
resulting guidance error grows quadratically, producing large miss
distances at engagement termination.

The contributions of this paper are threefold: (1)~a 2-D
point-mass engagement simulator with full EKF gating and PN
guidance; (2)~identification of a critical injection rate that
maximizes miss distance while maintaining sub-$5\%$ detection; and
(3)~Monte~Carlo validation with $N = 200$ engagements per
condition, supported by statistical significance testing and
power-law scaling analysis.  Section~\ref{sec:background} reviews
the analytical framework; Section~\ref{sec:simulation} describes
the simulator and baseline validation; Section~\ref{sec:results}
presents the critical-rate sweep and statistical results; and
Section~\ref{sec:discussion} discusses implications and
limitations.

%% ================================================================
\section{Background}
\label{sec:background}

\subsection{Proportional Navigation Guidance}
The lateral acceleration command of a PN-guided missile is
\begin{equation}
  a_{\text{cmd}} = N\, V_c\, \dot{\sigma}
  \label{eq:pn}
\end{equation}
where $N$ is the navigation ratio (typically $3$--$5$), $V_c$ is
the closing velocity, and $\dot{\sigma}$ is the line-of-sight
(LOS) angular rate.  Under ideal conditions, PN guarantees
collision when $N \ge 2$~\cite{zarchan2012}.

\subsection{Extended Kalman Filter with Innovation Gating}
The EKF tracks the target state
$\mathbf{x} = [x_t,\; y_t,\; \dot{x}_t,\; \dot{y}_t]^\top$
using range--bearing measurements.  The chi-squared innovation gate
rejects measurement $\mathbf{z}_k$ if the normalized innovation
\begin{equation}
  \gamma_k
  = \tilde{\mathbf{y}}_k^\top \mathbf{S}_k^{-1}
    \tilde{\mathbf{y}}_k
  > \chi^2_{\alpha,\,\nu}
  \label{eq:gate}
\end{equation}
exceeds the threshold, where
$\chi^2_{0.99,\,2} = 9.21$ for $99\%$ confidence with $\nu = 2$
degrees of freedom~\cite{barshalom2001}.  When a measurement is
rejected, the EKF propagates on its predicted state alone, causing
the covariance to grow.

\subsection{False Signal Injection Model}
Under PHANTOM injection, the corrupted bearing measurement is
\begin{equation}
  \beta_{\mathrm{meas}}(t)
  = \beta_{\mathrm{true}}(t) + \Delta\sigma(t) + \eta_\beta
  \label{eq:injection}
\end{equation}
where
$\Delta\sigma(t)
  = \int_0^{t} \delta\dot{\sigma}(\tau)\,d\tau$
is the accumulated false bearing angle and
$\eta_\beta \sim \mathcal{N}(0,\,\sigma_\beta^2)$ is sensor
noise.  For a linear ramp profile,
$\delta\dot{\sigma}(t) = r$ (constant), giving
$\Delta\sigma(t) = r\,(t - t_{\mathrm{start}})$.  The key insight
is that for sufficiently small~$r$, the instantaneous contribution
to $\gamma_k$ remains below the gate, so the EKF accepts the
corrupted measurement and its state estimate diverges from truth.

%% ================================================================
\section{Simulation Framework}
\label{sec:simulation}

\subsection{Engagement Geometry}
The simulator models a 2-D point-mass engagement in Cartesian
coordinates.  Table~\ref{tab:params} lists the nominal parameters
drawn from the PHANTOM configuration file.  Initial conditions
place the missile at the origin with the target at range $r_0$
along a collision course.  Target velocity is randomized within
$\pm 10\%$ of the nominal value across Monte~Carlo trials.

\begin{table}[t]
\caption{PHANTOM Simulation Parameters}
\label{tab:params}
\centering
\begin{tabular}{lll}
\toprule
Parameter & Value & Description \\
\midrule
$N$                       & 4.0       & Navigation ratio \\
$\Delta t$                & 0.02~s    & Update rate (50~Hz) \\
$\sigma_r$                & 10.0~m    & Range noise std \\
$\sigma_\beta$            & 0.00873~rad & Bearing noise std ($0.5^\circ$) \\
$\chi^2_{\mathrm{th}}$    & 9.21      & Gate threshold (99\%, 2~DOF) \\
$v_m$                     & 600~m/s   & Initial missile speed \\
$v_t$                     & 250~m/s   & Target speed \\
$r_0$                     & 5000~m    & Initial separation \\
\bottomrule
\end{tabular}
\end{table}

\subsection{Baseline Validation}
\label{subsec:baseline}
Before evaluating injection profiles, the simulator was validated
under no-injection conditions.  Table~\ref{tab:baseline}
summarizes the baseline results over $N = 100$ randomized
engagements with seed~42.  The mean miss distance of
$2.07 \pm 0.11$~m confirms correct PN guidance implementation
consistent with published benchmarks~\cite{zarchan2012}.
Fig.~\ref{fig:baseline} shows the miss-distance distribution.

\begin{table}[t]
\caption{Baseline Validation Results ($N = 100$, No Injection)}
\label{tab:baseline}
\centering
\begin{tabular}{lll}
\toprule
Metric & Value & Gate \\
\midrule
Mean miss distance  & $2.07 \pm 0.11$~m & $< 5$~m $\checkmark$ \\
95th percentile     & $2.29$~m           & $< 10$~m $\checkmark$ \\
Detection rate      & 0.00\%             & $= 0\%$ $\checkmark$ \\
Min / Max miss      & $1.86$ / $2.41$~m  & --- \\
\bottomrule
\end{tabular}
\end{table}

\begin{figure}[t]
  \centering
  \includegraphics[width=\columnwidth]
    {figures/phantom_fig01_baseline_miss_distribution}
  \caption{Baseline miss-distance distribution ($N = 100$,
  seed~42). Red dashed line marks the 5-m gate threshold.
  All engagements achieve sub-3-m miss, validating the PN
  guidance and EKF tracking implementation.}
  \label{fig:baseline}
\end{figure}

\subsection{Injection Profiles}
PHANTOM implements four injection profile types: (1)~\textbf{No
Injection} (baseline); (2)~\textbf{Linear Ramp}, a constant
false-LOS rate $r$ beginning at a configurable start time;
(3)~\textbf{Step}, an instantaneous bearing offset for testing
gate sensitivity; and (4)~\textbf{Sinusoidal}, an oscillating
injection for studying frequency-dependent detection.

The linear ramp profile was selected as the primary deception
mechanism because its slowly growing contribution to the innovation
statistic maximizes the time the EKF accepts corrupted measurements
before $\gamma_k$ exceeds the chi-squared gate.  All subsequent
results use this profile.

%% ================================================================
\section{Results}
\label{sec:results}

\subsection{Critical Rate Identification}
A sweep of 11 ramp rates from $0.005$ to $0.100$~rad/s was
conducted, each with $N = 200$ Monte~Carlo engagements
(Table~\ref{tab:sweep}).  The critical rate was defined as the
rate maximizing mean miss distance subject to a detection rate
below~$5\%$.  This was identified at $0.035$~rad/s, achieving a
mean miss of $20\,536 \pm 909$~m with $4.22\%$ detection
(Fig.~\ref{fig:critical_rate}).  The $95\%$ confidence interval
on the mean miss is $[20\,410,\; 20\,661]$~m.

\begin{table}[t]
\caption{Ramp Rate Sweep Results ($N = 200$ per rate)}
\label{tab:sweep}
\centering
\begin{tabular}{cccc}
\toprule
Rate (rad/s) & Mean Miss (m) & Std (m) & Detection (\%) \\
\midrule
0.005 & 20\,413 & 938 & 1.01 \\
0.010 & 20\,414 & 910 & 1.07 \\
0.015 & 20\,467 & 962 & 1.46 \\
0.020 & 20\,414 & 973 & 1.75 \\
0.025 & 20\,476 & 959 & 2.40 \\
0.030 & 20\,520 & 853 & 3.21 \\
0.032 & 20\,361 & 913 & 3.44 \\
\textbf{0.035} & \textbf{20\,536} & \textbf{909} & \textbf{4.22} \\
0.040 & 20\,345 & 895 & 7.28 \\
0.075 & 20\,414 & 886 & 57.96 \\
0.100 & 20\,319 & 1010 & 66.66 \\
\bottomrule
\multicolumn{4}{l}{\textbf{Bold}: Critical rate $0.035$~rad/s} \\
\end{tabular}
\end{table}

\begin{figure}[t]
  \centering
  \includegraphics[width=\columnwidth]
    {figures/phantom_fig05_critical_rate_combined}
  \caption{PHANTOM ramp rate optimization. Left: mean miss distance
  vs.\ injection rate with 200-m target (dashed). Right: EKF
  detection rate vs.\ injection rate with 5\% threshold (dashed).
  Vertical dashed line marks the critical rate of $0.035$~rad/s.}
  \label{fig:critical_rate}
\end{figure}

The bootstrap $95\%$ confidence interval on the critical rate
estimate is $[0.010,\; 0.035]$~rad/s (bootstrap std $= 0.0072$),
reflecting the flat miss-distance plateau across low rates.

\subsection{EKF Innovation Analysis}
Fig.~\ref{fig:innovation} shows the EKF innovation statistic
$\gamma_k$ over the course of a representative critical-rate
engagement (seed~392).  Below the gate threshold of $9.21$, all
measurements are accepted, and the EKF state estimate diverges
smoothly from the true target position.  Above the threshold,
measurements are rejected and the EKF propagates on its predicted
state.  The detection rate of $4.22\%$ indicates that the vast
majority of measurements pass the gate, confirming that the ramp
rate stays within the kinematically plausible envelope for most
of the engagement.

\begin{figure}[t]
  \centering
  \includegraphics[width=\columnwidth]
    {figures/phantom_fig06_innovation_timeline}
  \caption{EKF innovation statistic $\gamma_k$ during a
  critical-rate engagement ($r = 0.035$~rad/s). Red dashed
  line: chi-squared threshold $\chi^2_{0.99,2} = 9.21$.
  Shaded regions indicate timesteps where the measurement was
  rejected.}
  \label{fig:innovation}
\end{figure}

\subsection{Trajectory Deviation}
Fig.~\ref{fig:trajectory} compares a baseline engagement (left,
direct hit at $\sim 2$~m miss) with a PHANTOM injection engagement
(right, $> 20$~km miss).  Under injection, the corrupted EKF
estimate diverges from the true target position, causing the PN
guidance law to steer the missile away from the actual target.
The miss vector grows throughout the engagement.

\begin{figure*}[t]
  \centering
  \includegraphics[width=\textwidth]
    {figures/phantom_fig07_trajectory_comparison}
  \caption{Trajectory comparison. Left: baseline engagement
  (direct hit). Right: PHANTOM injection at $0.035$~rad/s
  (miss $> 20$~km). Blue: missile path; orange: target path;
  green dashed: EKF estimate.}
  \label{fig:trajectory}
\end{figure*}

\subsection{Statistical Analysis}
The difference in miss distance between baseline ($\bar{d} = 2.07$~m,
$N = 100$) and PHANTOM injection ($\bar{d} = 20\,536$~m, $N = 200$)
was evaluated using Welch's $t$-test.  The result was highly
significant: $t = 104.36$, $p < 0.001$, Cohen's $d = 58.16$,
corresponding to a mean increase of approximately $20\,534$~m.

Miss distance scales approximately as
$d_{\mathrm{miss}} \approx 21\,058 \cdot
  \dot{\sigma}_{\mathrm{ramp}}^{1.53}$
($R^2 = 0.940$) over the pre-saturation regime, consistent with
the quadratic false-angle accumulation predicted by the integral
in~\eqref{eq:injection} (Fig.~\ref{fig:scaling}).

\begin{figure}[t]
  \centering
  \includegraphics[width=\columnwidth]
    {figures/phantom_fig09_miss_distance_scaling}
  \caption{Power-law fit to miss distance vs.\ ramp rate over
  the pre-saturation envelope.  Exponent $b = 1.53$,
  $R^2 = 0.940$.}
  \label{fig:scaling}
\end{figure}

%% ================================================================
\section{Discussion}
\label{sec:discussion}

The PHANTOM results demonstrate that EKF innovation gating, while
effective at rejecting large instantaneous anomalies, is vulnerable
to slowly ramping false signals that accumulate over the
engagement timeline.  At the critical rate of $0.035$~rad/s, the
injection remains within the chi-squared acceptance envelope for
the majority of timesteps ($95.8\%$ acceptance rate), yet produces
miss distances exceeding $20$~km.  This suggests that the
chi-squared gate is a necessary but insufficient defense against
kinematically consistent false target injection.

Several limitations bound the present study.  The simulator uses
a 2-D point-mass model without aerodynamic drag, thrust dynamics,
or seeker gimbal limits.  A single missile type and a
non-maneuvering target are modeled.  No infrared counter-counter
measures (IRCCM) or adaptive gating are considered.  These
simplifications enable clean isolation of the EKF vulnerability
but limit direct operational applicability.

Phase~2 of the PHANTOM program will extend this framework in two
directions: (1)~an adaptive injection controller that modulates the
ramp rate in real time based on estimated $\gamma_k$ feedback, and
(2)~a large-language-model (LLM) policy trained on the Monte~Carlo
dataset to synthesize novel injection profiles.  The validated
Phase~1 simulator provides the deterministic, reproducible
environment required for training data generation.

%% ================================================================
\section{Conclusion}
\label{sec:conclusion}

This paper presented PHANTOM, a simulation framework for
kinematically consistent false LOS rate injection against
EKF-gated PN-guided missiles.  A linear ramp profile at
$0.035$~rad/s was shown to produce a mean miss distance of
$20\,536 \pm 909$~m with a $4.22\%$ detection rate across
$N = 200$ Monte~Carlo engagements.  The result is statistically
significant ($t = 104.36$, $p < 0.001$, $d = 58.16$) and
reproducible across random seeds.  Miss distance follows a
power-law scaling ($b = 1.53$, $R^2 = 0.940$) consistent with
the quadratic false-angle accumulation model.  The complete
codebase, data, and figures are archived at
\texttt{10.5281/zenodo.XXXXXX}.

\section*{Acknowledgment}
The author thanks the University of North Florida Department of
Aerospace and Mechanical Engineering for computational resources.

\begin{thebibliography}{99}

\bibitem{zarchan2012}
P.~Zarchan, \emph{Tactical and Strategic Missile Guidance},
6th~ed.\hskip 1em plus 0.5em minus 0.4em
Reston, VA: AIAA, 2012.

\bibitem{barshalom2001}
Y.~Bar-Shalom, X.~R.~Li, and T.~Kirubarajan,
\emph{Estimation with Applications to Tracking and Navigation}.
\hskip 1em plus 0.5em minus 0.4em
New York: Wiley, 2001.

\bibitem{siouris2004}
G.~M.~Siouris,
\emph{Missile Guidance and Control Systems}.
\hskip 1em plus 0.5em minus 0.4em
New York: Springer, 2004.

\bibitem{li2005}
X.~R.~Li and V.~P.~Jilkov,
``Survey of maneuvering target tracking---Part~V: Multiple-model
methods,''
\emph{IEEE Trans. Aerosp. Electron. Syst.}, vol.~41, no.~4,
pp.~1255--1321, Oct.~2005.

\bibitem{blair1998}
W.~D.~Blair and M.~Brandt-Pearce,
``Statistical description of monopulse parameters for tracking
Rayleigh targets,''
\emph{IEEE Trans. Aerosp. Electron. Syst.}, vol.~34, no.~2,
pp.~597--611, Apr.~1998.

\bibitem{fitzgerald1985}
R.~J.~Fitzgerald,
``Simple tracking filters: Closed-form solutions,''
\emph{IEEE Trans. Aerosp. Electron. Syst.}, vol.~AES-21, no.~6,
pp.~781--798, Nov.~1985.

\bibitem{skolnik2008}
M.~I.~Skolnik,
\emph{Radar Handbook}, 3rd~ed.\hskip 1em plus 0.5em minus 0.4em
New York: McGraw-Hill, 2008.

\bibitem{neri2001}
F.~Neri,
\emph{Introduction to Electronic Defense Systems}, 2nd~ed.
\hskip 1em plus 0.5em minus 0.4em
Raleigh, NC: SciTech, 2001.

\bibitem{adamy2015}
D.~L.~Adamy,
\emph{EW 104: Electronic Warfare Against a New Generation of
Threats}.
\hskip 1em plus 0.5em minus 0.4em
Norwood, MA: Artech House, 2015.

\bibitem{poisel2011}
R.~A.~Poisel,
\emph{Modern Communications Jamming: Principles and Techniques},
2nd~ed.\hskip 1em plus 0.5em minus 0.4em
Norwood, MA: Artech House, 2011.

\bibitem{brown2012}
R.~G.~Brown and P.~Y.~C.~Hwang,
\emph{Introduction to Random Signals and Applied Kalman Filtering},
4th~ed.\hskip 1em plus 0.5em minus 0.4em
New York: Wiley, 2012.

\bibitem{farooq2012}
A.~Farooq and D.~Mu,
``Adaptive sliding mode guidance of a homing missile,''
\emph{J. Chinese Inst. Eng.}, vol.~35, no.~3,
pp.~365--375, 2012.

\bibitem{gutman1998}
S.~Gutman and G.~Leitmann,
``Optimal strategies in the neighborhood of a collision course,''
\emph{AIAA J.}, vol.~14, no.~9, pp.~1210--1212, 1976.

\bibitem{shima2005}
T.~Shima and J.~Shinar,
``Time-varying linear pursuit-evasion game models with bounded
controls,''
\emph{J. Guid. Control Dyn.}, vol.~25, no.~3,
pp.~425--432, 2002.

\bibitem{yanushevsky2007}
R.~Yanushevsky,
\emph{Modern Missile Guidance}.
\hskip 1em plus 0.5em minus 0.4em
Boca Raton, FL: CRC Press, 2007.

\end{thebibliography}

\end{document}
